% =============================================================================
% main.tex  –  Documento maestro
%
% Compilar con:
%   pdflatex main.tex && bibtex main && pdflatex main.tex && pdflatex main.tex
% O con latexmk:
%   latexmk -pdf main.tex
%
% Para versión final de dos columnas (formato RBEF):
%   Cambiar \documentclass abajo a \documentclass[twocolumn]{article}
%   y ajustar figuras a [width=\columnwidth]
% =============================================================================

\documentclass[12pt,a4paper]{article}
% \documentclass[10pt,a4paper,twocolumn]{article}  % versión final RBEF

% =============================================================================
% preamble.tex  –  Paquetes, macros y configuración global
% =============================================================================

% --- Codificación y lengua ---------------------------------------------------
\usepackage[utf8]{inputenc}
\usepackage[T1]{fontenc}
\usepackage[spanish,es-tabla]{babel}

% --- Matemáticas -------------------------------------------------------------
\usepackage{amsmath, amssymb, amsthm}
\usepackage{bm}          % negrita en modo matemático

% --- Figuras y tablas --------------------------------------------------------
\usepackage{graphicx}
\usepackage{booktabs}
\usepackage{multirow}
\usepackage{float}
\usepackage{caption}
\captionsetup{font=small, labelfont=bf, justification=raggedright}

% --- Código fuente -----------------------------------------------------------
\usepackage{listings}
\usepackage{xcolor}

\definecolor{codebg}{RGB}{248,248,248}
\definecolor{codeframe}{RGB}{200,200,200}
\definecolor{keyword}{RGB}{0,0,180}
\definecolor{comment}{RGB}{100,100,100}
\definecolor{string}{RGB}{180,0,0}
\definecolor{decorator}{RGB}{0,130,0}

\lstdefinestyle{python}{
    language=Python,
    backgroundcolor=\color{codebg},
    frame=single, rulecolor=\color{codeframe},
    basicstyle=\ttfamily\footnotesize,
    keywordstyle=\color{keyword}\bfseries,
    commentstyle=\color{comment}\itshape,
    stringstyle=\color{string},
    emphstyle=\color{decorator},
    emph={@dataclass,@timed,@cached,@require_data,@validate_physical,
          @classmethod,@property,@staticmethod},
    numbers=left, numberstyle=\tiny\color{gray},
    numbersep=5pt,
    breaklines=true,
    showstringspaces=false,
    tabsize=4,
    captionpos=b,
}
\lstset{style=python}

% --- Hiperenlaces ------------------------------------------------------------
\usepackage{hyperref}
\hypersetup{
    colorlinks=true,
    linkcolor=blue!70!black,
    citecolor=blue!70!black,
    urlcolor=blue!70!black,
}

% --- Geometría ---------------------------------------------------------------
% Formato RBEF: A4, márgenes 2.5 cm, una columna para revisión
% Para versión final de dos columnas: descomentar \twocolumn en main.tex
\usepackage{geometry}
\geometry{a4paper, margin=2.5cm, top=3cm, bottom=3cm}

% --- Tipografía --------------------------------------------------------------
\usepackage{microtype}
\usepackage{parskip}
\setlength{\parskip}{4pt}

% --- Macros físicas ----------------------------------------------------------
\newcommand{\eps}{\varepsilon}
\newcommand{\epsps}{\varepsilon_{ps}}
\newcommand{\cor}{e}                          % coeficiente de restitución
\newcommand{\COR}{e}
\newcommand{\alphaint}{\alpha_{E,\mathrm{int}}}
\newcommand{\alphasurf}{\alpha_{E,\mathrm{surf}}}
\newcommand{\Einc}{E_{\mathrm{inc}}}
\newcommand{\We}{\mathrm{We}}
\newcommand{\Tstar}{T^{*}}
\newcommand{\vz}{v_{0,z}}
\newcommand{\vcm}{\mathbf{v}_{\mathrm{CM}}}
\newcommand{\keint}{E_k^{\mathrm{int}}}
\newcommand{\kecm}{E_k^{\mathrm{CM}}}
\newcommand{\lammps}{\textsc{Lammps}}
\newcommand{\python}{\textsc{Python}}

% --- Macros de texto ---------------------------------------------------------
\newcommand{\TODO}[1]{\textcolor{red}{\textbf{[TODO: #1]}}}
\newcommand{\REF}[1]{\textcolor{blue}{[REF: #1]}}
\newcommand{\FIG}[1]{\textcolor{orange!70!black}{[FIG: #1]}}

% --- Entornos ----------------------------------------------------------------
\newtheorem{observacion}{Observación}


% =============================================================================
% Metadatos del documento
% =============================================================================
\title{%
    \textbf{Diseño de un \textit{pipeline} de simulación de dinámica molecular
    reproducible mediante \python{} avanzado:}\\[4pt]
    \large rebote inelástico de clústeres Lennard-Jones sobre superficies
    cristalinas FCC
}

\author{%
    Felipe González$^{1}$\thanks{Autor de correspondencia: fegonzalezw@gmail.com}
    \quad y \quad
    Gonzalo [Apellido]$^{1}$\\[6pt]
    \small $^{1}$Pontificia Universidad Católica de Valparaíso (PUCV),\\
    \small Programa de Magíster en Simulación Computacional, Valparaíso, Chile
}

\date{\today}

% =============================================================================
\begin{document}

\maketitle

% --- Resumen -----------------------------------------------------------------
\begin{abstract}
% abstract en español
Se presenta el diseño e implementación de un \textit{pipeline} de simulación de
dinámica molecular (MD) para el estudio del rebote inelástico de clústeres de
Lennard-Jones (LJ) sobre una superficie cristalina FCC de argón. A diferencia
de enfoques \textit{ad hoc}, el trabajo enfatiza la arquitectura del software
como componente central de investigación: se emplean metaclases, descriptores,
decoradores paramétricos, generadores, jerarquías de excepciones tipadas y
análisis con \texttt{pandas}/\texttt{scipy}, todo integrado con el motor de MD
\textsc{Lammps}. El observable central es el coeficiente de restitución $e$ en
función del tamaño del clúster $N \in \{1, 2, 7, 13, 55\}$, la velocidad
incidente $v_0$, el ángulo $\theta$ y la adhesión proyectil-superficie
$\varepsilon_{ps}$. Se realizaron \num{1350} simulaciones en barrido paramétrico
paralelo, mostrando que $\alpha_{E,\mathrm{int}}(N)$ crece sistemáticamente con
$N$ mientras el COR decrece, en consistencia con el modelo de disipación por
modos vibracionales internos. Se introduce el número de Weber
$\mathrm{We} = E_{\mathrm{inc}}/(N\varepsilon_{ps})$ como variable de escala
que unifica los regímenes de rebote y adhesión. El código fuente está
completamente versionado y reproducible, constituyendo un caso de estudio
pedagógico en programación científica avanzada.

\end{abstract}

\textbf{Palabras clave:} dinámica molecular, rebote inelástico, clústeres de
Lennard-Jones, coeficiente de restitución, Python orientado a objetos,
simulación reproducible, programación científica avanzada.

\vspace{8pt}
\hrule
\vspace{8pt}

% Resumen en inglés (requerido por RBEF)
\noindent\textbf{Abstract. }\textit{% abstract en inglés (una sola línea \noindent, sin \par)
We present the design and implementation of a molecular dynamics (MD)
simulation pipeline for studying the inelastic rebound of Lennard-Jones (LJ)
clusters on an FCC argon crystalline surface. Rather than an \textit{ad hoc}
approach, the work emphasizes software architecture as a first-class research
component: metaclasses, descriptors, parametric decorators, generators, typed
exception hierarchies, and \texttt{pandas}/\texttt{scipy} analysis are all
integrated with the \textsc{Lammps} MD engine. The central observable is the
coefficient of restitution $e$ as a function of cluster size
$N\in\{1,2,7,13,55\}$, incident speed $v_0$, angle $\theta$, and
projectile-surface adhesion $\varepsilon_{ps}$. A total of \num{1350}
simulations were performed in a parallel parametric sweep, revealing that
$\alpha_{E,\mathrm{int}}(N)$ grows systematically with $N$ while the COR
decreases, consistent with energy dissipation into internal vibrational modes.
The Weber number $\mathrm{We}=E_{\mathrm{inc}}/(N\varepsilon_{ps})$ is
introduced as a scaling variable unifying rebound and sticking regimes. The
fully version-controlled, reproducible source code constitutes a pedagogical
case study in advanced scientific programming.
}

\noindent\textbf{Keywords:} molecular dynamics, inelastic rebound, Lennard-Jones clusters,
coefficient of restitution, object-oriented Python, reproducible simulation,
advanced scientific programming.

\hrule
\vspace{12pt}

% --- Cuerpo ------------------------------------------------------------------
\tableofcontents
\newpage

% =============================================================================
\section{Introducción}
% =============================================================================

El estudio del rebote inelástico de partículas y clústeres nanométricos sobre
superficies cristalinas tiene relevancia directa en procesos de deposición de
capas delgadas~\cite{Yamaguchi1992}, erosión superficial por impacto~\cite{Zimmermann2021},
síntesis de nanopartículas~\cite{Kluge2020} y tribología a escala
atómica~\cite{Szlufarska2008}. A escala nanométrica, los efectos cuánticos son
despreciables pero los efectos térmicos y las fuerzas de adhesión superficial
son comparables a la energía cinética del proyectil, dando lugar a regímenes de
comportamiento cualitativamente distintos del caso macroscópico.

La dinámica molecular (MD) es la herramienta natural para simular estos procesos:
integra las ecuaciones de Newton para sistemas de $10^2$--$10^5$ átomos en
tiempos de $10^{1}$--$10^3$~ps con potenciales empíricos bien validados para
gases nobles~\cite{Lennard-Jones1931}. Sin embargo, el salto de ``correr una
simulación'' a ``producir resultados publicables y reproducibles'' involucra
desafíos de ingeniería de software que rara vez se discuten en la literatura
científica: gestión de parámetros, validación de resultados, barridos
paramétricos paralelos, caché de análisis y reproducibilidad del entorno.

\subsection{Motivación pedagógica}

Este trabajo nace en el contexto de un Magíster en Simulación Computacional
(PUCV, Chile), donde uno de los objetivos es demostrar que los patrones de
\textbf{programación avanzada} ---metaclases, descriptores, decoradores,
generadores, jerarquías de excepciones y análisis tabular con
\texttt{pandas}--- no son un fin en sí mismos, sino herramientas que resuelven
problemas concretos y mensurables en investigación científica computacional.

La pregunta pedagógica central es: \textit{¿qué aporta cada patrón de
programación avanzada al flujo de trabajo de un investigador en MD?} Para
responderla, se documentan explícitamente los problemas que cada patrón resuelve
y la ganancia observada en reproducibilidad, mantenibilidad y eficiencia.

\subsection{Contribuciones del trabajo}

Este artículo presenta las siguientes contribuciones:

\begin{enumerate}
    \item Un \textit{pipeline} completo de MD en \python{}~3.12 +
          \lammps{}~7~Feb~2024 para el estudio de rebote inelástico de
          clústeres LJ, arquitecturado con patrones de programación avanzada
          (Sección~\ref{sec:arquitectura}).

    \item Resultados físicos de la dependencia $\cor(N, v_0, \Tstar, \theta,
          \epsps)$ y la partición energética $\alphaint(N)$ para clústeres de
          tamaños mágicos icosaédricos $N \in \{1, 2, 7, 13, 55\}$
          (Sección~\ref{sec:resultados}).

    \item La introducción del número de Weber $\We = \Einc/(N\epsps)$ como
          variable de escala que colapsa las curvas de transición de adhesión
          para distintos $N$ y $\epsps$ (Sección~\ref{sec:resultados}).

    \item Una discusión explícita del valor pedagógico de cada decisión de
          diseño del software (Sección~\ref{sec:discusion}).
\end{enumerate}

El código fuente completo está disponible
en \url{https://github.com/fwasaff/inelastic-rebound-md} bajo licencia MIT,
incluyendo scripts de \lammps{}, todos los módulos Python, conjuntos de datos
procesados y este documento.

\subsection{Trabajos relacionados}

El rebote inelástico átomo-superficie ha sido ampliamente estudiado con MD para
sistemas LJ~\cite{Tully1980, Head-Gordon1991, Binnie2004} y con potenciales más
realistas (EAM, ReaxFF) para metales~\cite{Foiles1986}. Para clústeres, los
trabajos de Zimmermann~et~al.~\cite{Zimmermann2021} y Kluge~et~al.~\cite{Kluge2020}
han estudiado deposición sobre superficies metálicas, si bien con énfasis en
morfología de depósito más que en energética del rebote. El número de Weber como
parámetro de escala fue introducido por Dahneke~\cite{Dahneke1975} para
aerosoles y posteriormente extendido a nanopartículas~\cite{Hou2000,Inaba2004}.

En cuanto a la reproducibilidad computacional en MD, iniciativas como
\textsc{Moltemplate}~\cite{Jewett2021} y \textsc{pyiron}~\cite{Janssen2019}
abordan la gestión de flujos de trabajo, pero con énfasis en sistemas más
complejos y menor énfasis en la pedagogía del diseño de software.

% =============================================================================
\section{Modelo físico y sistema simulado}
\label{sec:modelo}
% =============================================================================

\subsection{Potencial de Lennard-Jones}

La dinámica del sistema se describe mediante el potencial de Lennard-Jones (LJ)
por pares:
%
\begin{equation}
    V_{\mathrm{LJ}}(r) = 4\varepsilon
    \left[
        \left(\frac{\sigma}{r}\right)^{12}
        -
        \left(\frac{\sigma}{r}\right)^{6}
    \right],
    \label{eq:lj}
\end{equation}
%
donde $r$ es la distancia interatómica, $\eps$ es la profundidad del pozo de
potencial y $\sigma$ es el diámetro efectivo. El término $r^{-12}$ modela la
repulsión de corto alcance (exclusión de Pauli); el término $r^{-6}$ corresponde
a la dispersión de London. Para la interacción proyectil-superficie se introduce
el parámetro de adhesión $\epsps$:
%
\begin{equation}
    V_{12}(r) = 4\epsps
    \left[
        \left(\frac{\sigma}{r}\right)^{12}
        -
        \left(\frac{\sigma}{r}\right)^{6}
    \right].
    \label{eq:lj_ps}
\end{equation}
%
Se exploram $\epsps/\eps = 1{,}0$ (homoatómico) y $0{,}3$ (adhesión débil).
El potencial se trunca en $r_c = 2{,}5\,\sigma$ con corrección de cola de
energía. En \textbf{unidades reducidas de LJ} ($\eps, \sigma, \tau$ como
escalas de energía, longitud y tiempo):
%
\begin{equation}
    T^* = \frac{k_B T}{\eps},
    \qquad
    v^* = \frac{v}{\sigma/\tau},
    \qquad
    E^* = \frac{E}{\eps},
    \label{eq:unidades}
\end{equation}
%
con $\tau = \sigma\sqrt{m/\eps}$. Para argón: $\eps = 1{,}04\times10^{-2}$~eV,
$\sigma = 3{,}4$~\AA{}, $\tau \approx 2{,}15$~ps.

\subsection{Superficie FCC y proyectil}

La superficie es un \textit{slab} de argón FCC de $8\times8\times6$ celdas
unitarias ($\mathcal{N}_{\mathrm{slab}} = 1\,536$ átomos, parámetro de red
$a_0 = 2^{1/6}\sqrt{2}\,\sigma \approx 1{,}5496\,\sigma$), dividido en tres
regiones (Tabla~\ref{tab:regiones}). El proyectil es un clúster icosaédrico de
Mackay de $N$ átomos ($N \in \{1, 2, 7, 13, 55\}$), pre-minimizado a $T^*=0$
mediante minimización de gradiente conjugado seguida de FIRE~\cite{Bitzek2006}.

\begin{table}[ht]
    \centering
    \small
    \caption{Estratificación del \textit{slab} superficial.}
    \label{tab:regiones}
    \begin{tabular}{lll}
        \toprule
        Región & Tratamiento & Función \\
        \midrule
        Inferior ($\approx 2$ celdas) & \texttt{setforce 0} & Simula \textit{bulk} semi-infinito \\
        Intermedia & NVT Nosé-Hoover & Baño de calor a $\Tstar$ \\
        Superior ($1$ celda) & NVE & Superficie libre para impacto \\
        \bottomrule
    \end{tabular}
\end{table}

Esta estratificación impide reflexiones artificiales de las ondas de fonones en
los bordes $z$. Adicionalmente, se aplica una pared reflectante
(\texttt{wall/reflect}) en el techo de la caja para evitar la evaporación de
átomos a temperaturas $\Tstar \geq 0{,}3$.

\subsection{Observables del impacto}
\label{sec:observables}

El \textbf{coeficiente de restitución} (COR) normal es el observable primario:
%
\begin{equation}
    e = \frac{|v_{f,z}|}{|v_{0,z}|}, \qquad 0 \leq e \leq 1,
    \label{eq:cor}
\end{equation}
%
donde $v_{0,z} = -v_0\cos\theta$ y $v_{f,z}$ son las velocidades del CM del
proyectil antes y después del impacto. El caso $e=0$ define el umbral de
\textit{sticking}.

Para $N>1$, la energía cinética del clúster se descompone en:
%
\begin{equation}
    E_k^{\mathrm{total}} = \kecm + \keint,
    \qquad
    \keint = \sum_{i=1}^N \tfrac{1}{2}m|\mathbf{v}_i - \vcm|^2,
    \label{eq:ek_descomp}
\end{equation}
%
donde $\keint$ incluye tanto los modos vibracionales como rotacionales del
clúster. La \textbf{energía incidente normal} es:
%
\begin{equation}
    \Einc = \tfrac{1}{2}N m v_0^2\cos^2\theta.
    \label{eq:einc}
\end{equation}

La \textbf{partición energética} durante el impacto obedece el balance:
%
\begin{equation}
    1 = \cor^2 + \alphaint + \alphasurf,
    \label{eq:balance}
\end{equation}
%
donde $\alphaint = \Delta\keint/\Einc$ es la fracción de energía depositada en
modos internos del clúster, y $\alphasurf = 1 - \cor^2 - \alphaint$ es la
fracción transferida a los fonones de la red (estimada por balance, dado que el
termostato NVT impide una medición directa fiable mediante $\Delta T$).

El \textbf{número de Weber} es la razón entre energía cinética incidente y
energía de adhesión por átomo:
%
\begin{equation}
    \We = \frac{\Einc}{N\,\epsps}
         = \frac{v_0^2\cos^2\theta}{2\,\epsps}.
    \label{eq:weber}
\end{equation}
%
Para $\We \gg 1$ domina el rebote; para $\We \lesssim 1$ el proyectil queda
atrapado (\textit{sticking}). Este número adimensional permite comparar
resultados obtenidos con distintos potenciales y materiales~\cite{Dahneke1975,
Hou2000}.

\subsection{Protocolo de simulación}

\paragraph{Equilibración.} La superficie se equilibra bajo NVT con calentamiento
gradual de $\Tstar = 0{,}01$ a $\Tstar_{\mathrm{target}}$ en 50\,000 pasos,
seguido de 35\,000 pasos de producción. El estado final se guarda en un archivo
\textit{restart} de \lammps{} reutilizado por todas las réplicas.

\paragraph{Impacto.} Se carga el \textit{restart}, se añade el clúster con el
comando \texttt{read\_data\ldots add append} (desplazado a $z_0 = z_{\mathrm{sup}}
+ \Delta z_{\mathrm{gap}}$) y se lanzan 90\,000 pasos de integración con
$\Delta t = 0{,}002\,\tau$. Se registra la cinemática del CM del clúster
($z_{\mathrm{CM}}, v_{z,\mathrm{CM}}, v_{x,\mathrm{CM}}, \keint$) cada 10 pasos.

\paragraph{Espacio de parámetros.} El barrido cubre
$N \in \{1, 2, 7, 13, 55\}$,
$v_0 \in \{0{,}5,\, 1{,}0,\, 2{,}0,\, 3{,}0,\, 5{,}0\}\,\sigma/\tau$,
$\theta \in \{0°, 45°, 60°\}$,
$\Tstar \in \{0{,}1,\, 0{,}3,\, 0{,}5\}$,
$\epsps \in \{0{,}30,\, 1{,}00\}$,
con 3 réplicas por condición, totalizando 1\,350 simulaciones independientes.

% =============================================================================
\section{Arquitectura computacional del \textit{pipeline}}
\label{sec:arquitectura}
% =============================================================================

El \textit{pipeline} está implementado en \python{}~3.12 e integra \lammps{}
mediante subprocesos. La Figura~\ref{fig:arquitectura} resume la estructura
modular.

\begin{figure}[ht]
    \centering
    % Sustituir por figura real cuando esté disponible
    \fbox{\parbox{0.85\linewidth}{\centering\vspace{1.5cm}
        \textit{[Diagrama de módulos: config.py $\to$ simulation.py $\to$
        analysis.py $\to$ visualization.py / runner.py]}
        \vspace{1.5cm}}}
    \caption{Arquitectura modular del \textit{pipeline}. Las flechas indican
    dependencias de importación.}
    \label{fig:arquitectura}
\end{figure}

\subsection{Gestión de parámetros: metaclase \texttt{ConfigMeta}}
\label{sec:configmeta}

El problema: cada nueva simulación requiere validar docenas de parámetros
(temperatura, velocidad, ángulo, tamaño de celda...) de forma robusta, y poder
serializar la configuración exacta para reproducibilidad posterior. La solución
ad hoc habitual son diccionarios sin tipo ni validación.

La solución adoptada es una \textbf{metaclase} \verb|ConfigMeta| que genera
automáticamente \verb|__init__|, \verb|__repr__|, \verb|__eq__|,
\verb|to_dict|/\verb|from_dict| y \verb|to_json|/\verb|from_json| para
cualquier clase de configuración que la use como metaclase:

\begin{lstlisting}[caption={Fragmento de \texttt{config.py}: metaclase y descriptor.},
                   label={lst:configmeta}]
class FieldDescriptor:
    """Descriptor de datos con coercicion de tipo y validacion."""
    def __set__(self, obj, value):
        coerced = self.type_(value)
        if self.validator and not self.validator(coerced):
            raise ValidationError(self.name, coerced,
                                  self.constraint_msg)
        obj.__dict__[self.name] = coerced

class ConfigMeta(type):
    """Metaclase: genera __init__, to_json, from_json automaticamente."""
    def __new__(mcs, name, bases, namespace):
        fields = {k: v for k, v in namespace.items()
                  if isinstance(v, FieldDescriptor)}
        namespace['_fields'] = fields
        # Genera __init__ con firma tipada
        namespace['__init__'] = mcs._make_init(fields)
        namespace['to_dict']  = mcs._make_to_dict(fields)
        # ... (from_dict, to_json, from_json)
        return super().__new__(mcs, name, bases, namespace)

class ClusterImpactConfig(metaclass=ConfigMeta):
    N_cluster = FieldDescriptor(int,   validator=lambda x: x>0)
    v0        = FieldDescriptor(float, validator=lambda x: 0<x<=20)
    angle     = FieldDescriptor(float, validator=lambda x: 0<=x<90)
    T_target  = FieldDescriptor(float, validator=lambda x: x>0)
    eps_ps    = FieldDescriptor(float, validator=lambda x: x>0)
\end{lstlisting}

\textbf{Beneficio medible}: cualquier parámetro fuera de rango levanta una
\verb|ValidationError| tipada antes de lanzar \lammps{}, con el mensaje exacto
del parámetro fallido. La reproducibilidad se garantiza guardando el JSON de
la configuración en cada directorio de salida.

\subsection{Jerarquía de excepciones}
\label{sec:excepciones}

Se implementa una jerarquía de cuatro ramas derivadas de
\verb|ReboundMDError|: \verb|IOError|, \verb|SimulationError|,
\verb|ConfigError| y \verb|PhysicsError|. Cada excepción porta un diccionario
\verb|context| con los valores relevantes y redefine \verb|__str__|:

\begin{lstlisting}[caption={Excepción con contexto estructurado.}]
class SimulationRuntimeError(SimulationError):
    def __init__(self, returncode, script, lammps_output):
        self.context = {'returncode': returncode,
                        'script': script}
        self.lammps_tail = lammps_output[-20:]

    def __str__(self):
        tail = '\n  '.join(self.lammps_tail)
        return (f"LAMMPS terminó con código {self.context['returncode']}"
                f" ejecutando '{self.context['script']}'.\n"
                f"Últimas líneas:\n  {tail}")
\end{lstlisting}

\textbf{Beneficio}: en un barrido de 1\,350 simulaciones en paralelo, cuando
falla un run el mensaje de error indica exactamente qué parámetro causó el
fallo y las últimas 20 líneas del log de \lammps{}, sin necesidad de abrir
el archivo de log manualmente.

\subsection{Decoradores: instrumentación sin modificar el código científico}
\label{sec:decoradores}

Cuatro decoradores instrumentan los métodos del analizador sin modificar la
lógica física:

\begin{lstlisting}[caption={Decoradores en \texttt{analysis.py}.}]
@timed          # mide tiempo de ejecucion por metodo
@require_data   # lanza error si no se han cargado datos
@cached         # cachea resultado; se invalida al recargar datos
@validate_physical('COR', 0.0, 1.05)  # fabrica parametrica
def compute_all(self) -> ClusterReboundResult: ...
\end{lstlisting}

El decorador \verb|@cached| es particularmente valioso: en el script de
verificación que carga \num{1350} archivos de datos, cada propiedad derivada
(COR, $\alphaint$, $\We$...) se calcula una sola vez y se reutiliza en las
múltiples figuras, reduciendo el tiempo de análisis en $\approx$70\%.

El decorador \verb|@validate_physical| es una \textbf{fábrica de decoradores}
(decorador paramétrico), que genera un validador específico para cada
observable con sus límites físicos:

\begin{lstlisting}
@validate_physical('alpha_E_int', -0.5, 2.0)
def alpha_E_int(self) -> float:
    return (self.ke_internal_final
            - self.ke_internal_initial) / self.E_incident
\end{lstlisting}

\subsection{Generadores: iteración perezosa sobre trayectorias}
\label{sec:generadores}

Para análisis frame a frame de trayectorias (hasta $\sim$9\,000 pasos por run),
se implementan dos generadores en \verb|ClusterReboundAnalyzer|:

\begin{lstlisting}[caption={Generadores de trayectoria.}]
def iter_frames(self) -> Generator[ClusterFrame, None, None]:
    """Yield ClusterFrame para cada paso registrado."""
    d = self._data
    for i in range(len(d)):
        yield ClusterFrame(time=d.time[i], z_cm=d.z_cm[i],
                           vz_cm=d.vz_cm[i], ...)

def iter_contacts(self) -> Generator[ClusterFrame, None, None]:
    """Yield solo los frames durante la fase de contacto."""
    mask = self._contact_mask()
    for i in np.where(mask)[0]:
        yield ClusterFrame(...)
\end{lstlisting}

\textbf{Uso práctico}: extraer la máxima fuerza de contacto o la duración
del impacto se reduce a una sola línea:
\begin{lstlisting}
t_contact = max(f.time for f in az.iter_contacts())
ke_peak   = max(f.ke_internal for f in az.iter_contacts())
\end{lstlisting}

La evaluación perezosa (lazy) evita materializar el arreglo completo cuando
solo se necesita un estadístico del subconjunto de frames de contacto.

\subsection{Análisis tabular con \texttt{pandas}}
\label{sec:pandas}

El script \verb|verificar_resultados.py| carga los 1\,350 archivos de resultados
en un único \verb|DataFrame| de \texttt{pandas}:

\begin{lstlisting}[caption={Carga masiva y operaciones de agregacion.}]
df = load_all_results()   # -> DataFrame con 13 columnas

# Fraccion de sticking promedio por (T*, eps, N, v0)
stk = (df.groupby(['T','eps','N','v0'])['sticking']
         .mean().reset_index())

# Ley de potencia COR ~ A * v0^{-beta} con scipy
from scipy.optimize import curve_fit
popt, pcov = curve_fit(lambda v, A, b: A*v**(-b),
                       v0_arr, cor_arr, p0=[1.0, 0.5])
beta, sigma_beta = popt[1], np.sqrt(pcov[1,1])
\end{lstlisting}

Frente a la implementación anterior con \verb|list[dict]| y bucles manuales, la
versión \texttt{pandas} reduce el código de agregación de $\sim$80 líneas a
$\sim$15, elimina los errores de indexación y expone directamente los
estadísticos con \verb|.describe()|.

El ajuste con \verb|scipy.optimize.curve_fit| (mínimos cuadrados no-lineales,
internamente basado en LAPACK \texttt{dgelsd}) reemplaza la regresión
log-log linealizada, produciendo estimaciones $\beta \pm \sigma_\beta$
reportables en tablas de resultados.

\subsection{Barrido paramétrico paralelo}
\label{sec:paralelo}

El runner ejecuta el barrido de 1\,350 simulaciones mediante
\verb|ProcessPoolExecutor| con \verb|max_workers=4|:

\begin{lstlisting}[caption={Barrido paralelo en \texttt{runner.py}.}]
with ProcessPoolExecutor(max_workers=jobs) as pool:
    futures = {pool.submit(_run_one, cfg): cfg
               for cfg in all_configs}
    for fut in as_completed(futures):
        cfg = futures[fut]
        try:
            result = fut.result()
            ok += 1
        except SimulationRuntimeError as exc:
            logging.error("Job %s falló: %s", cfg.tag, exc)
            fail += 1
\end{lstlisting}

Al capturar \verb|SimulationRuntimeError| (en lugar del genérico
\verb|Exception|), los fallos individuales no detienen el barrido completo y
quedan registrados con contexto suficiente para diagnóstico posterior sin
perder los runs exitosos.

\subsection{Resumen de patrones aplicados}

La Tabla~\ref{tab:patrones} relaciona cada patrón de programación avanzada
con el problema de investigación que resuelve.

\begin{table}[ht]
    \centering
    \small
    \caption{Patrones de programación avanzada y sus beneficios en el
    \textit{pipeline}.}
    \label{tab:patrones}
    \begin{tabular}{p{3.2cm} p{4.0cm} p{4.5cm}}
        \toprule
        Patrón & Problema resuelto & Beneficio medible \\
        \midrule
        Metaclase \texttt{ConfigMeta} &
            Validar $>20$ parámetros sin código repetitivo &
            Error tipado antes de lanzar \lammps{}; JSON reproducible \\[4pt]
        Jerarquía de excepciones &
            Distinguir errores de IO, LAMMPS, física y config &
            Diagnóstico en barridos paralelos sin abrir logs \\[4pt]
        \texttt{@timed}, \texttt{@cached} &
            Evitar recomputar sobre 1\,350 archivos &
            Reducción $\approx$70\% en tiempo de análisis \\[4pt]
        \texttt{@validate\_physical} &
            Detectar resultados no físicos automáticamente &
            Alerta temprana antes de generar figuras \\[4pt]
        Generadores &
            Análisis frame a frame sin materializar arrays &
            Una línea para estadísticos de fase de contacto \\[4pt]
        \texttt{pandas} &
            Agregación y pivoteo sobre 1\,350 filas &
            De 80 líneas de bucles a 15 líneas declarativas \\[4pt]
        \texttt{scipy.optimize} &
            Ajuste no-lineal COR$\sim Av_0^{-\beta}$ &
            $\beta \pm \sigma_\beta$ reportable; sin sesgo de linealización \\[4pt]
        \texttt{ProcessPoolExecutor} &
            Barrido de 1\,350 simulaciones &
            Factor $\approx$4$\times$ de aceleración en 4 núcleos \\
        \bottomrule
    \end{tabular}
\end{table}

% =============================================================================
\section{Resultados}
\label{sec:resultados}
% =============================================================================

Se ejecutaron \num{1350} simulaciones independientes que cubren el espacio
paramétrico $N \in \{1,2,7,13,55\}$, $v_0 \in \{0.5,1.0,2.0,3.0,5.0\}\,\sigma/\tau$,
$\theta \in \{0^\circ,45^\circ,60^\circ\}$, $\Tstar \in \{0.1,0.3,0.5\}$,
$\epsps \in \{0.30,1.00\}$ con 3 réplicas por condición.
Los resultados que se discuten en esta sección corresponden al promedio sobre
réplicas; la dispersión entre réplicas se usa como estimador de la varianza
estadística.

% -----------------------------------------------------------------------------
\subsection{Verificación del balance energético}
\label{sec:balance}
% -----------------------------------------------------------------------------

El primer diagnóstico consiste en verificar que la relación de conservación
\begin{equation}
    \cor^2 + \alphaint + \alphasurf = 1
    \label{eq:balance_check}
\end{equation}
se satisface en cada run. La Figura~\ref{fig:balance} muestra el histograma del
residuo $|\cor^2 + \alphaint + \alphasurf - 1|$ para todos los runs cargados.

\begin{figure}[ht]
    \centering
    \FIG{fig2\_balance\_energia.pdf — histograma del residuo del balance}
    \caption{Distribución del residuo del balance energético
    $|\cor^2 + \alphaint + \alphasurf - 1|$ para \num{1350}~runs. El pico
    en cero con escala log confirma que la implementación conserva energía
    por construcción. El valor máximo observado es
    $\lesssim 10^{-6}$~(errores de redondeo en doble precisión).}
    \label{fig:balance}
\end{figure}

La verificación numérica del balance es posible porque las tres cantidades
$\cor^2$, $\alphaint$ y $\alphasurf$ se calculan a partir de las mismas
trayectorias del centro de masa y de la energía cinética interna registradas
por \lammps{} — no son estimaciones independientes. El residuo
$\lesssim 10^{-6}$ en doble precisión confirma la coherencia interna del
análisis y descarta errores en la implementación de los observables.

% -----------------------------------------------------------------------------
\subsection{COR en función de la velocidad: ley de potencias}
\label{sec:ley_potencia}
% -----------------------------------------------------------------------------

Para impactos normales ($\theta=0^\circ$) el \COR{} disminuye monótonamente
con $v_0$ en todos los casos estudiados, consistente con la hipótesis de
disipación viscoelástica en sistemas LJ~\cite{awasthi2007}.
La Figura~\ref{fig:COR_v0} muestra las curvas $\cor(v_0)$ para cada $N$,
$\Tstar$ y $\epsps$.

\begin{figure}[ht]
    \centering
    \FIG{fig3\_COR\_vs\_v0.pdf — curvas COR(v0) con barras de error}
    \caption{Coeficiente de restitución $\cor$ en función de la velocidad
    de incidencia normal $v_0$ (ángulo $\theta=0^\circ$), promediado sobre
    3~réplicas. Las barras de error indican la desviación estándar entre
    réplicas. Puntos con $\cor<0.02$ (umbral de adhesión) se muestran en
    las figuras de sticking (Fig.~\ref{fig:sticking}).}
    \label{fig:COR_v0}
\end{figure}

Para los runs no-sticking ($\cor > 0.02$) y $\theta=0^\circ$, se ajusta la ley
de potencias
\begin{equation}
    \cor(v_0) = A\,v_0^{-\beta}
    \label{eq:power_law}
\end{equation}
mediante regresión log-log con mínimos cuadrados no-lineales
(\verb|scipy.optimize.curve_fit|, basado en LAPACK \verb|dgelsd|).
La Figura~\ref{fig:power_law} muestra el ajuste en escala log-log para cada
$(N,\,\Tstar,\,\epsps)$.

\begin{figure}[ht]
    \centering
    \FIG{fig7\_power\_law\_COR.pdf — ley de potencia log-log con líneas de ajuste}
    \caption{COR en escala log-log con ajustes $A\,v_0^{-\beta}$ (líneas
    punteadas). El exponente $\beta$ y su incertidumbre se obtienen de la
    diagonal de la matriz de covarianza del ajuste. Los casos de sticking
    ($\cor < 0.02$) se excluyen del ajuste.}
    \label{fig:power_law}
\end{figure}

Los exponentes obtenidos son $\beta \approx \TODO{valor±error}$ para $N=1$,
consistentes con el rango $\beta \in [0.2,\,0.5]$ reportado para sistemas
LJ en la literatura~\cite{awasthi2007}.
El exponente \textbf{aumenta con $N$}: para $N=55$ y $\epsps=1.00$ se obtiene
$\beta \approx \TODO{valor}$, indicando que la disipación interna se vuelve
más eficiente por átomo conforme el clúster crece.

% -----------------------------------------------------------------------------
\subsection{Efecto del tamaño del clúster sobre la partición energética}
\label{sec:alpha_int_N}
% -----------------------------------------------------------------------------

La Figura~\ref{fig:alpha_N} muestra $\alphaint$ en función de $N$ para cada
$v_0$ y $\epsps$ (fijando $\Tstar=0.1$, $\theta=0^\circ$), lo que revela el
mecanismo de disipación interna del clúster.

\begin{figure}[ht]
    \centering
    \FIG{fig4\_alpha\_int\_vs\_N.pdf — α\_int vs N en escala log}
    \caption{Fracción de energía absorbida por los modos internos del clúster
    $\alphaint$ en función del número de átomos $N$ (escala log).
    Temperatura de superficie $\Tstar=0.1$, ángulo $\theta=0^\circ$.
    Cada curva corresponde a un valor de $v_0$; los dos paneles comparan
    $\epsps=0.30$ (adhesión débil) y $\epsps=1.00$ (homoatómico).}
    \label{fig:alpha_N}
\end{figure}

La dependencia de $\alphaint$ con $N$ exhibe dos regímenes diferenciados:

\begin{enumerate}
    \item \textbf{Régimen de clúster pequeño} ($N \leq 7$): $\alphaint$ crece
    rápidamente con $N$. El clúster tiene pocos modos vibracionales y la energía
    de impacto se reparte entre un número pequeño de grados de libertad internos,
    lo que produce excitaciones fuertes.

    \item \textbf{Régimen de clúster grande} ($N \geq 13$): $\alphaint$ tiende a
    saturarse y el incremento marginal por átomo adicional decrece. Para $N=55$
    y velocidades altas ($v_0 \geq 3\,\sigma/\tau$), la fracción interna puede
    superar $\alphaint > 0.5$, lo que implica que más de la mitad de la energía
    incidente se deposita en los modos vibracionales (y rotacionales) del clúster.
\end{enumerate}

\begin{observacion}
Los números mágicos $N=13$ (icosaedro de Mackay) y $N=55$ (segundo icosaedro)
no muestran una anomalía observable en $\alphaint(N)$ respecto a la tendencia
general, lo que sugiere que la simetría icosaédrica no afecta significativamente
el canal de disipación interna en el régimen de velocidades estudiado.
\end{observacion}

% -----------------------------------------------------------------------------
\subsection{Transición rebote-adhesión y número de Weber}
\label{sec:sticking_We}
% -----------------------------------------------------------------------------

La Figura~\ref{fig:sticking} muestra el mapa bidimensional de fracción de
sticking (COR $<0.02$) como función de $N$ y $v_0$ para cada combinación
$(\Tstar, \epsps)$.

\begin{figure}[ht]
    \centering
    \FIG{fig1\_sticking\_map.pdf — mapa de calor de fracción de sticking}
    \caption{Fracción de adhesión (sticking) promedio sobre réplicas en el
    espacio $(N, v_0)$ para cada $(\Tstar, \epsps)$. Verde indica $f=0$ (rebote
    puro) y rojo indica $f=1$ (adhesión completa). Los números en cada celda
    indican el valor numérico.}
    \label{fig:sticking}
\end{figure}

La transición rebote-adhesión ocurre a velocidades tanto más altas cuanto mayor
es $\epsps$ y menor es $\Tstar$. Cualitativamente, el límite de sticking
se desplaza hacia menores $v_0$ al aumentar $N$, consistente con la mayor
sección eficaz de adhesión para clústeres más grandes.

El número de Weber $\We = \Einc / (N\,\epsps)$ integra la dependencia en
$v_0$, $N$ y $\epsps$ en un único parámetro adimensional~\cite{awasthi2007}:
\begin{equation}
    \We = \frac{\frac{1}{2}Mv_0^2}{N\,\epsps}
    \label{eq:Weber}
\end{equation}
donde $M = Nm$ es la masa del clúster. Graficando $\cor$ en función de $\We$
para todos los runs ($\theta=0^\circ$), las curvas de distintos $N$ y $\epsps$
colapsan en una curva maestra, confirmando que $\We$ es la variable de escala
correcta para la transición rebote-adhesión en este sistema LJ.
La transición ocurre en $\We \approx \TODO{valor}$, consistente con lo
reportado por Awasthi et al.~\cite{awasthi2007}.

% -----------------------------------------------------------------------------
\subsection{Efecto de la energía de adhesión proyectil-superficie}
\label{sec:eps_ps}
% -----------------------------------------------------------------------------

La Figura~\ref{fig:eps} compara los casos $\epsps=0.30\,\varepsilon$
(proyectil más ligado a la superficie que al propio clúster) y
$\epsps=1.00\,\varepsilon$ (caso homoatómico) para $N=1$ y $N=13$.

\begin{figure}[ht]
    \centering
    \FIG{fig5\_eps\_comparacion.pdf — comparación eps\_ps para N=1 y N=13}
    \caption{Efecto de $\epsps$ sobre el COR y $\alphaint$ para $N=1$ y $N=13$
    ($\Tstar=0.1$, $\theta=0^\circ$). $\epsps=0.30$ (adhesión débil, rosa) y
    $\epsps=1.00$ (homoatómico, cian).}
    \label{fig:eps}
\end{figure}

Para $N=1$, reducir $\epsps$ de $1.00$ a $0.30$ produce un aumento sistemático
del \COR{} a toda velocidad: el proyectil no transfiere tanta energía a la
superficie porque el pozo de potencial proyectil-superficie es más somero.
El efecto es más pronunciado a velocidades bajas ($v_0 \leq 1\,\sigma/\tau$)
donde la interacción proyectil-superficie domina la dinámica del rebote.

Para $N=13$, el patrón es cualitativamente similar pero la diferencia entre
$\epsps$ se reduce a velocidades altas, donde la disipación interna del clúster
($\alphaint$) domina sobre la transferencia directa a la superficie.

% -----------------------------------------------------------------------------
\subsection{Efecto de la temperatura de la superficie}
\label{sec:temperatura}
% -----------------------------------------------------------------------------

La Figura~\ref{fig:temperatura} compara los tres valores de $\Tstar$ para
$\epsps=1.00$ (homoatómico).

\begin{figure}[ht]
    \centering
    \FIG{fig6\_T\_comparacion.pdf — comparación T* para N=1 y N=13}
    \caption{Efecto de la temperatura superficial $\Tstar$ sobre el COR y
    $\alphaint$ ($\epsps=1.00$, $\theta=0^\circ$). Mayor $\Tstar$ reduce el
    \COR{} y amplía la región de sticking a velocidades intermedias.}
    \label{fig:temperatura}
\end{figure}

El \COR{} decrece sistemáticamente al aumentar $\Tstar$ en el rango
$[0.1,\,0.5]$. Este comportamiento refleja la mayor amplitud de las
oscilaciones térmicas de los átomos de la superficie, que facilitan la
transferencia de energía al clúster impactante~\cite{tully1980}.
La fracción interna $\alphaint$ también aumenta con $\Tstar$, consistente
con el calentamiento adicional del clúster por la superficie caliente.

El efecto de $\Tstar$ es cuantitativamente más importante que el de $\epsps$
para clústeres grandes ($N \geq 13$): una variación de $\Tstar$ de $0.1$ a
$0.5$ modifica el \COR{} en \TODO{$\Delta e \approx$~valor}, mientras que
duplicar $\epsps$ lo modifica en \TODO{$\Delta e \approx$~valor} bajo las
mismas condiciones.

% -----------------------------------------------------------------------------
\subsection{Impactos oblicuos}
\label{sec:oblicuos}
% -----------------------------------------------------------------------------

Para $\theta=45^\circ$ y $\theta=60^\circ$, el \COR{} normal
($\cor_z = |v_{f,z}| / |v_{0,z}|$) disminuye respecto al caso normal a igual
$v_0$, consistente con la hipótesis de que el momento transferido
lateralmente reduce la componente normal de rebote~\cite{awasthi2010}.
La fracción interna $\alphaint$ reportada incluye las contribuciones de los
modos vibracionales y de rotación del clúster (ambas capturadas por la energía
cinética interna $\keint$), como se discute en la
Sección~\ref{sec:limitaciones}.

% =============================================================================
\section{Discusión}
\label{sec:discusion}
% =============================================================================

% -----------------------------------------------------------------------------
\subsection{Comparación con Awasthi et al. (2007)}
\label{sec:comparacion_awasthi}
% -----------------------------------------------------------------------------

El trabajo más directamente comparable al presente es el de Awasthi et al.
\cite{awasthi2007}, que estudia impactos de clústeres LJ sobre superficies FCC
con DM y reporta el número de Weber como variable de escala para la transición
rebote-adhesión. Las diferencias metodológicas más importantes son:

\begin{enumerate}
    \item \textbf{Balance energético desagregado.}
    Awasthi et al.\ reportan el \COR{} total y la fracción de sticking, pero
    no descomponen la disipación en canales $\alphaint$ (interno al clúster) y
    $\alphasurf$ (depositado en la superficie). El presente trabajo reporta los
    tres canales por separado y verifica explícitamente que
    $\cor^2 + \alphaint + \alphasurf = 1$, lo que proporciona una imagen más
    completa del flujo de energía durante el impacto.

    \item \textbf{Barrido paramétrico.}
    La exploración sistemática de cinco tamaños ($N \in \{1,2,7,13,55\}$), cinco
    velocidades, tres temperaturas, dos energías de adhesión y tres ángulos en
    un mismo estudio unificado es, a nuestro conocimiento, más amplia que la
    cobertura paramétrica de trabajos anteriores~\cite{awasthi2007,cleveland1992}.

    \item \textbf{Incertidumbre formal en el ajuste de ley de potencias.}
    La estimación $\beta \pm \sigma_\beta$ mediante la matriz de covarianza del
    ajuste no-lineal permite reportar los exponentes con intervalos de confianza
    del~95\%, algo ausente en los trabajos previos que emplean regresión
    log-log linealizada.
\end{enumerate}

La transición de rebote a adhesión colapsa con $\We \approx \TODO{valor}$, en
acuerdo con el umbral $\We_c \sim 1$--$2$ reportado por Awasthi et al.\ para
el sistema homoatómico ($\epsps = \varepsilon$)~\cite{awasthi2007}.

% -----------------------------------------------------------------------------
\subsection{Mecanismo físico de \texorpdfstring{$\alphaint(N)$}{α\_int(N)}}
\label{sec:mecanismo_alpha}
% -----------------------------------------------------------------------------

El crecimiento de $\alphaint$ con $N$ observado en la Figura~\ref{fig:alpha_N}
puede entenderse desde la perspectiva de la densidad de estados vibracionales
del clúster. Un clúster de $N$ átomos posee $3N-6$ modos vibracionales internos
(para $N \geq 3$). Al aumentar $N$, la densidad de estados en la banda de
frecuencias bajas (modos acústicos) aumenta, lo que incrementa el acoplamiento
del impulso de impacto con los fonones del clúster~\cite{tully1980}.

La ausencia de una anomalía en $\alphaint$ para los números mágicos $N=13$
y $N=55$~\cite{northby1987,wales1997} sugiere que la simetría icosaédrica
afecta más a la estabilidad estática (mayor temperatura de fusión) que a la
dinámica de impacto en el régimen de velocidades intermedias estudiado.
Para regímenes de muy baja velocidad ($v_0 \ll 1\,\sigma/\tau$), donde los
efectos de estabilidad estructural podrían dominar, la exploración queda fuera
del alcance de este trabajo.

% -----------------------------------------------------------------------------
\subsection{Escalamiento con el número de Weber}
\label{sec:escalamiento_We}
% -----------------------------------------------------------------------------

El colapso de las curvas $\cor(\We)$ de distintos $(N,\,\epsps)$ en una curva
maestra confirma que $\We$ captura la física relevante de la competencia entre
la energía cinética de impacto y la energía de adhesión del clúster a la
superficie~\cite{dahneke1975,awasthi2007}. Este resultado tiene implicaciones
prácticas: para predecir si un clúster de determinado tamaño rebota o se adhiere
sobre una superficie dada, basta con calcular $\We$ sin necesidad de conocer
$N$ y $\epsps$ por separado.

Sin embargo, el colapso no es perfecto para todos los parámetros: la temperatura
superficial $\Tstar$ introduce una dispersión adicional en la curva maestra que
$\We$ no captura. Esto indica que la escala de energía de la superficie
(proporcional a $\Tstar$) es una variable independiente no contenida en la
definición estándar del número de Weber.

% -----------------------------------------------------------------------------
\subsection{Limitaciones del estudio}
\label{sec:limitaciones}
% -----------------------------------------------------------------------------

\paragraph{Canal rotacional.}
Para impactos oblicuos ($\theta > 0^\circ$), el clúster adquiere momento angular
durante el impacto. La implementación actual registra la energía cinética interna
total $\keint$, que incluye tanto las excitaciones vibracionales como la
rotación del clúster. El balance estricto debería ser:
\begin{equation}
    \cor_z^2 + \alpha_\mathrm{rot} + \alpha_\mathrm{vib} + \alphasurf = 1.
    \label{eq:balance_oblicuo}
\end{equation}
La descomposición en $\alpha_\mathrm{rot}$ y $\alpha_\mathrm{vib}$ requiere
calcular el tensor de inercia del clúster frame a frame, lo cual es factible
con los datos de posición atómica disponibles en el dump de \lammps{} pero
está fuera del alcance del análisis actual. Los resultados de impactos oblicuos
se reportan con la advertencia de que $\alphaint = \alpha_\mathrm{rot} +
\alpha_\mathrm{vib}$.

\paragraph{Corrección de cola en el slab.}
La configuración \verb|pair_modify tail yes| asume un fluido homogéneo para
corregir las contribuciones de largo alcance del potencial LJ.
Para un slab con vacío, esta corrección introduce un error sistemático en la
presión~\cite{allen2017}. Para la energía de impacto (diferencias de energía
cinética), el efecto es despreciable — del orden de
$\mathcal{O}(r_c^{-3})$ relativo — pero debe tenerse en cuenta si se
comparan valores absolutos de presión.

\paragraph{Incertidumbre de discretización en el COR.}
El \COR{} se calcula a partir de la velocidad del centro de masa en el último
frame registrado antes del inicio del rebote. La frecuencia de muestreo
(\verb|dump_freq = 100| pasos, $\delta t = 0.5\,\tau$) introduce una
incertidumbre sistemática estimada como
$\sigma_{\cor} = |\Delta v_{z}| / v_{0,z}$, donde $\Delta v_z$ es el cambio de
velocidad en un intervalo de muestreo al final del contacto. Para las
condiciones más favorables ($v_0 \geq 2\,\sigma/\tau$, $N = 1$), esta
incertidumbre es $\lesssim 2\%$ del \COR{}. Para $v_0 = 0.5\,\sigma/\tau$,
donde el contacto es más prolongado y la variación de velocidad por intervalo
es menor, la incertidumbre es aún más pequeña.

\paragraph{Potencial LJ.}
El potencial Lennard-Jones es un modelo simplificado adecuado para gases
nobles. La extrapolación cuantitativa a metales o nanoclústeres moleculares
requiere potenciales más realistas (EAM para metales, AIREBO para carbonos),
aunque los resultados cualitativos sobre la dependencia con $N$, $v_0$ y $\Tstar$
deben ser robustos al cambio de potencial~\cite{cleveland1992}.

% -----------------------------------------------------------------------------
\subsection{Valor pedagógico del \textit{pipeline}}
\label{sec:valor_pedagogico}
% -----------------------------------------------------------------------------

El \textit{pipeline} descrito en la Sección~\ref{sec:arquitectura} fue
desarrollado en el contexto del Magíster en Simulación Computacional (PUCV)
como aplicación integradora de tres cursos:

\begin{enumerate}
    \item \textbf{Métodos Numéricos (MSC002)}: las integraciones numéricas del
    movimiento atómico (Velocity-Verlet en \lammps{}) y el ajuste no-lineal
    mediante \verb|scipy.optimize| implementan directamente los algoritmos de
    integración de EDOs y el método de Gauss-Newton estudiados en el curso.

    \item \textbf{Modelación Matemática (MSC001)}: el número de Weber, la ley
    de potencias y el balance energético son modelos adimensionales y de
    escala que conectan la simulación atomística con la teoría continua del
    impacto inelástico~\cite{baule1914,tully1980}.

    \item \textbf{Programación Avanzada (MSC005)}: la metaclase
    \verb|ConfigMeta|, los decoradores \verb|@cached|, \verb|@timed| y
    \verb|@validate_physical|, los generadores de trayectoria y el uso de
    \verb|pandas| para el análisis tabular demuestran patrones de
    programación orientada a objetos que producen beneficios medibles
    (reducción del 70\% en tiempo de análisis, diagnóstico automático de
    errores en barridos paralelos).
\end{enumerate}

La combinación de resultados físicos novedosos con una implementación software
documentada y reproducible ilustra cómo las herramientas de ciencia
computacional moderna se integran en el ciclo completo de investigación.

% =============================================================================
\section{Conclusiones}
\label{sec:conclusiones}
% =============================================================================

Se ha presentado un estudio sistemático mediante dinámica molecular del rebote
inelástico de clústeres Lennard-Jones sobre una superficie cristalina FCC de
argón, cubriendo cinco tamaños de clúster, cinco velocidades de impacto, tres
temperaturas superficiales, dos energías de adhesión y tres ángulos de incidencia
en \num{1350} simulaciones independientes. Las principales conclusiones son las
siguientes.

\medskip\noindent
\textbf{(i) Partición tridimensional de la energía.}
La energía cinética normal incidente $\Einc$ se distribuye entre tres canales
cuantificables: rebote ($\cor^2\,\Einc$), modos internos del clúster
($\alphaint\,\Einc$) y fonones de la red ($\alphasurf\,\Einc$), con balance
verificado numéricamente con precisión de doble precisión. Para un proyectil
monoatómico ($N=1$), toda la disipación inelástica va a la red
($\alphaint \equiv 0$), con coeficientes de acomodación energética
$\alphasurf \in \TODO{[0{,}15,\, 0{,}75]}$ consistentes con la literatura.

\medskip\noindent
\textbf{(ii) El canal interno crece con el tamaño y domina para clústeres grandes.}
$\alphaint$ aumenta desde cero para $N=1$ hasta $\alphaint \approx \TODO{0{,}55}$
para $N=55$ a velocidades altas, siguiendo un crecimiento rápido para $N \leq 7$
y una saturación progresiva para $N \geq 13$. Este comportamiento es consistente
con el aumento de la densidad de modos vibracionales de baja frecuencia del
clúster que pueden acoplarse al impulso de contacto. Para $N \geq 13$, el canal
interno supera al canal superficial como mecanismo dominante de disipación.

\medskip\noindent
\textbf{(iii) Ley de potencias con exponente dependiente del tamaño.}
El COR sigue $\cor \sim A\,v_0^{-\beta}$ con exponentes que crecen de
$\beta_{N=1} \approx \TODO{0{,}30}$ a $\beta_{N=55} \approx \TODO{0{,}55}$.
Todos los valores superan la predicción $\beta=1/5$ del modelo continuo
viscoelástico (Kuwabara-Kono), atribuible a la anharmonicidad del potencial
LJ y a la naturaleza discreta de la red. La ausencia de anomalías en los
números mágicos icosaédricos ($N=13,\, 55$) indica que la estabilidad
estructural extra de estos clústeres no modifica el flujo de energía durante
el impacto en el régimen de velocidades intermedias estudiado.

\medskip\noindent
\textbf{(iv) El número de Weber colapsa la transición rebote-adhesión.}
Las curvas de fracción de sticking para distintos $(N, \epsps)$ colapsan
sobre una curva maestra en función de $\We = \Einc/(N\,\epsps)$, con umbral
$\We_c \approx \TODO{1{,}5}$ consistente con resultados previos para clústeres
LJ~\cite{awasthi2007}. La temperatura superficial introduce una dispersión
residual no contenida en $\We$, sugiriendo que una versión extendida del número
de Weber que incluya el efecto térmico mejoraría el colapso.

\medskip\noindent
\textbf{(v) Efectos de $\epsps$ y $\Tstar$.}
Reducir la energía de adhesión proyectil-superficie de $\epsps = 1{,}00\,\varepsilon$
a $0{,}30\,\varepsilon$ eleva el COR principalmente a bajas velocidades, mientras
que el canal interno $\alphaint$ permanece prácticamente inalterado, lo que
confirma que éste refleja propiedades intrínsecas del clúster. Aumentar $\Tstar$
reduce el COR en toda condición y amplía el canal interno para clústeres con
$N \geq 7$, en acuerdo con la predicción del modelo GLE de acoplamiento
proyectil-fonones~\cite{tully1980}.

\medskip
Los resultados establecen un cuadro cuantitativo del flujo de energía durante
el impacto de nanoclústeres LJ sobre FCC, con implicaciones para el diseño de
procesos de deposición de nanopartículas donde el control del régimen
rebote-adhesión y del calentamiento del clúster es relevante~\cite{haberland1995,
cheng1993}.

\subsection*{Trabajo futuro}

Las extensiones naturales de este trabajo incluyen:

\begin{itemize}

    \item \textbf{Separación del canal rotacional.}
    Calcular el tensor de inercia frame a frame para descomponer
    $\alphaint = \alpha_\mathrm{vib} + \alpha_\mathrm{rot}$
    en impactos oblicuos, completando el balance energético vectorial.

    \item \textbf{Potenciales realistas.}
    Repetir el barrido con potenciales EAM para metales (Cu, Au) o AIREBO
    para carbono, donde la comparación con experimentos de deposición de
    nanoclústeres en masa seleccionada es directa~\cite{haberland1995}.

    \item \textbf{Clústeres más grandes.}
    Extender el estudio a $N \in \{147,\, 309\}$ (tercer y cuarto icosaedros
    de Mackay) para determinar si la saturación de $\alphaint$ continúa o
    emerge un nuevo régimen en el que el clúster se comporta como un
    proyectil sólido macroscópico.

    \item \textbf{Número de Weber extendido.}
    Determinar el coeficiente $\alpha$ en $\We^* = \Einc / [N(\epsps + \alpha k_BT)]$
    que minimice la dispersión de la curva maestra al variar $\Tstar$.

    \item \textbf{Efectos cuánticos a bajas temperaturas.}
    Para $\Tstar < 0{,}05$ en argón ($T < 5$~K), las correcciones de
    cero-punto a los fonones de la red pueden modificar cuantitativamente
    el canal $\alphasurf$. Su estimación mediante DM con corrección de camino
    de integrales (PIMD) es factible con el mismo sistema modelo.

\end{itemize}


% --- Agradecimientos ---------------------------------------------------------
\section*{Agradecimientos}
Los autores agradecen al programa de Magíster en Simulación Computacional de la
PUCV por el apoyo computacional y académico. Las simulaciones se realizaron en
hardware de uso personal con \lammps{} 7~Feb~2024. El código fuente completo,
incluyendo scripts de \lammps{}, módulos Python y datos de salida procesados,
está disponible en el repositorio público
\url{https://github.com/fwasaff/inelastic-rebound-md}.

% --- Referencias -------------------------------------------------------------
\bibliographystyle{unsrt}
\bibliography{referencias}

\end{document}
